\documentclass[a4paper, 11pt]{article}
\usepackage[utf8]{inputenc}
\usepackage[T1]{fontenc}
\usepackage[margin=1.5cm]{geometry}

\title{Praca inżynierska - możliwe źródłą/odwołania}
\author{Szymon Wasążnik\\325522}
\date{}

\begin{document}
\maketitle

\begin{itemize}
  \setlength{\itemsep}{0pt}
  \setlength{\parskip}{0pt}
  % algorytmy gen
  \item D. E. Goldberg "Algorytmy genetyczne i ich zastosowanie" Wydawnictwo Naukowo-Techniczne,\\Warszawa 1998 - 2003
  \item A.E. Eiben, J.E. Smith "Introduction to Evolutionary Computing", Springer Berlin, Heidelberg, 2015
  \item Jarosław Arabas "Wykłady z algorytmów ewolucyjnych", Wydawnictwo Naukowo-Techniczne,\\Warszawa, 2001 - 2004
  \item J. Cytowski "Algorytmy genetyczny. Podstawy i zastosowania", Akademicka Oficyna Wydawnicza PLJ, Warszawa, 1996
  % sztuczna
  \item "Sztuczna inteligencja dla inżynierów. Metody ogólne", Oficyna Wydawnicza Politechniki Warszawskiej, Warszawa, 2022
  \item L. Rutkowski "Metody i techniki sztucznej inteligencji", PWN, Warszawa, 2009
  % python
  \item Allen B. Downey "Think Python: How to Think Like a Computer Scientist" 2nd Edition, 2016
  \item Al Sweigart "Beyond the Basic Stuff with Python: Best Practices for Writing Clean Code", 2020
  \item Charles Severance "Python for Everybody: Exploring Data in Python 3", 2016 % TODO: check if has any use
  \item Wes McKinney "Python for Data Analysis: Data Wrangling with pandas, NumPy, and Jupyter" 3rd Edition, 2022
  \item "Introduction to Machine Learning with Python: A Guide for Data Scientists"
  % to check if manual selection of architecture can be better
  \item Francois Chollet "Deep Learning with Python, Second Edition" 2nd Edition, 2021
\end{itemize}
\end{document}